\section{Discussion}
\label{sec:discussion}

Honest patient networks, built only from what was known at the time, are organised around clinical outcomes well beyond chance. They are highly clustered, split into strong communities, and those communities differ sharply in their outcome rates. Leakage behaves very differently for the two questions. For long stay, the outcome is a length of time rather than a diagnosis, so the future record adds detail but not the answer. For new high blood pressure, the future record contains the diagnosis itself, and the cheating network becomes more aligned with the outcome. Even then the effect looks modest in summary numbers, a reminder that leakage can hide inside small changes.

Simple network measures also revealed three problems in the data before any model was trained: the 2015 billing-code switch, which splits both kinds of networks when raw codes are used; connections between stays of the same patient, which carry much of the high-blood-pressure signal; and some official high-blood-pressure cases whose patients already had the diagnosis. Each could quietly inflate the results of a prediction model.

The diagnosis families agree with medical knowledge: breathing failure groups with pneumonia and pleural effusion, sepsis with kidney failure, and heart failure with cardiomyopathy. Pleural effusion, the strongest long-stay giveaway in the future record, belongs to the breathing-failure family, consistent with long stays being driven by serious lung and heart illness.

\paragraph{Limitations.} This study describes patterns and does not test predictions on new patients. We use one way of building the networks (diagnoses only, 10 nearest neighbours); adding medications or lab tests, or changing the number of neighbours, could change the picture. Homophily is computed with all outcomes known, which is fine for description but would itself be leakage in a prediction setting. Our rebuilt labels differ slightly from Stanford's official ones. Finally, a connection in the diagnosis network means two diagnoses tend to occur within a year of each other, not that one causes the other, and our grouping ignores which came first.
